\documentclass[11pt]{article}
\usepackage[margin=1in]{geometry}
\usepackage[T1]{fontenc}
\usepackage{lmodern,microtype,amsmath,amssymb,amsthm,mathtools}
\usepackage{booktabs,longtable,array,xurl,hyperref,listings}
\hypersetup{hidelinks,pdftitle={Unknot recognition: executable baseline and specification audit},
 pdfauthor={Implementation report}}
\lstset{basicstyle=\small\ttfamily,columns=fullflexible,breaklines=true,
 keepspaces=true,showstringspaces=false,frame=single}
\newtheorem{theorem}{Theorem}[section]
\newtheorem{proposition}[theorem]{Proposition}
\newtheorem{lemma}[theorem]{Lemma}
\theoremstyle{remark}\newtheorem{remark}[theorem]{Remark}
\newcommand{\F}{\mathbb F_2}
\newcommand{\Q}{\mathbb Q}
\newcommand{\Z}{\mathbb Z}
\newcommand{\Kh}{\widetilde{\operatorname{Kh}}}
\newcommand{\rank}{\operatorname{rank}}
\newcommand{\code}[1]{\texttt{\detokenize{#1}}}
\title{Unknot recognition:\\an executable baseline and an audit of the\\quasi-polynomial specification}
\author{Implementation and technical report}
\date{17 September 2026}
\begin{document}
\maketitle
\begin{center}
\fbox{\begin{minipage}{0.92\linewidth}
\textbf{Delivery status: partial implementation.}\quad
The requested general $n^{O(\log n)}$ implementation has \emph{not} been
established. The supplied general recognizer is executable and mathematically
exact, but its complete backend has an exponential $2^{O(n)}$ bound.
The hierarchy components described below do not constitute a complete
implementation of the algorithm announced in the supplied notes.
\end{minipage}}
\end{center}
\begin{abstract}
This archive contains a dependency-free Python reference recognizer, three
implemented hierarchy-related components, tests, examples, and a complexity
audit. The complete recognition path constructs reduced Khovanov homology over
$\F$, with optional descending-diagram and determinant certificates before it.
The independent hierarchy components implement compressed cohomology-dual
normal coordinates, essential boundary-pattern testing on a known 3-ball, and
bounded lexicographic accounting. We specify their mathematical contracts,
explain what is verified, and identify the geometric operations and quantitative
invariants still needed for the requested running time. Tests and small-case
measurements are reported as such, not as a substitute for a worst-case proof.
\end{abstract}

\section{The target and the source audit}
The input is a classical one-component knot diagram with $n$ crossings. The
requested result is a two-valued decision procedure, correct on every such
input, with worst-case bit complexity
\[
 \exp\bigl(O((\log(n+2))^2)\bigr)=n^{O(\log n)}
\]
for $n\ge2$. Returning \code{unknown} on difficult inputs does not meet that
specification. Nor does an exponential algorithm become quasi-polynomial
because its tested examples happen to finish quickly.

The supplied February 2021 slides announce this bound on pages 11--14
\cite{slides}. Their starting hierarchy algorithm is summarized on page 53.
Pages 66--74 identify four accelerations: a quadratic surface-complexity bound,
efficient hierarchy encoding, multisurfaces, and generalized Heegaard splittings
with Cheeger regions. The last page names the substantial geometric procedures
required to implement that strategy. The present work follows the distinction
in the slides between the initial hierarchy method and those additional
accelerations.

An additional source located during this work is Lackenby's July 2026 preprint
\cite{lackenby2026}. Section 9 explicitly distinguishes iteration counting from
running-time analysis and leaves some operations insufficiently specified for
the latter. Section 10 supplies a linear hierarchy-length bound in its own
complexity measure, not the logarithmic bound needed here. This is a statement
about that particular source, not a claim that the announced result is false
or that a complete implementation is impossible.

The missing implementation is concrete: this archive does not construct the
required hierarchical multisurfaces, perform compressed surgery with pattern
history, find and use Cheeger regions, or execute the weak reductions in the
final flowchart. It therefore supplies neither the entire geometric algorithm
nor the combined quantitative proof. Instead of replacing those operations by
oracles, the archive separates a complete, slower recognizer from real,
individually testable hierarchy components.

\section{Executable interface and input validation}
The package is \code{unknotlab}. It requires Python 3.10 or newer and only the
standard library. It runs directly from the extracted folder:
\begin{lstlisting}
python -m unknotlab examples/torus_3_5.json --check-d2
python -m unittest discover -s tests -v
python benchmark.py
python hierarchy_demo.py
\end{lstlisting}
The CLI accepts either a planar-diagram object or a braid-closure object:
\begin{lstlisting}
{"pd": [[1,2,3,4], [4,3,5,6], [6,5,2,1]]}
{"braid": {"strands": 3, "word": [1,-2,1,-2]}}
\end{lstlisting}
A crossing $(a,b,c,d)$ lists outgoing half-edges in counterclockwise cyclic
order. The underpass joins $a$ to $c$ and the overpass joins $b$ to $d$.
Every positive integer arc label must occur exactly twice. Labels are normalized
before use. Empty PD means one crossing-free circle; hidden additional
crossing-free components are not part of this encoding.

Let $\alpha$ pair the two darts carrying the same arc label, and let $\sigma$
rotate the four darts of each crossing. The cycles of $\sigma\alpha$ are the
faces of the orientable ribbon surface. The parser first checks that joining
opposite darts and arc-paired darts gives exactly one knot component. This
also makes the projection graph connected. The equation
\[
 V-E+F=n-2n+F=2
\]
then verifies that its rotation system is spherical. Incidence data that
encode a knot only on a positive-genus supporting surface are rejected.

A braid generator $i>0$ has the strand at position $i$ over position $i+1$;
negative generators reverse that crossing. The word is read from top to
bottom and its endpoints are closed in their original positions. Component
count is checked from the permutation. The parser rejects a strand count
larger than the word length plus one before allocating a permutation: that
many strands cannot be connected by the available crossings.

The ordinary output statuses are \code{unknot} and \code{knotted}; both have
exit code zero. A selected resource limit or an inconclusive
\code{--fast-only} run gives \code{unknown}, exit code two. Invalid input gives
exit code three. Neither timeout nor resource exhaustion is interpreted as a
mathematical certificate. No resource limit is imposed by default, so the
complete procedure may consume impractically large resources.

Complexity expressed just in $n$ assumes the conventional arc labels
$1,\dots,2n$. For arbitrary larger labels or irrelevant extra input text, one
must also account for the actual input bit length. This parsing cost is not
hidden in the crossing count.

\section{The complete recognition path}
\subsection{Polynomial-time sufficient certificates}
The first optional test tries every starting dart in both traversal directions.
When a traversal meets every crossing first on its overpass, it is descending.
Such a diagram is an unknot diagram \cite{descending}: after choosing the
basepoint, its open arc can be realized with strictly decreasing height, with
a closing arc outside the rest of the diagram. The classical straightening of
this configuration gives an unknotted circle. The test examines $O(n)$ starts
and $O(n)$ crossing encounters for each. Failure of this test says nothing.

The second test joins semiarcs through overpasses to obtain the $n$ Wirtinger
arcs. It forms the integral coloring matrix with a row
\[
 2x_{\mathrm{over}}-x_{\mathrm{under},1}-x_{\mathrm{under},2}=0
\]
for every crossing, adding coefficients when arcs coincide. The absolute
value of an $(n-1)$-square cofactor is the knot determinant
$|\Delta_K(-1)|$ \cite{kolay}. The code uses fraction-free Bareiss elimination,
with row pivoting and a check of every exact division. A determinant other
than one certifies nontriviality because the unknot has determinant one.
\emph{Determinant one is never used as a positive certificate.}

\subsection{The explicitly constructed reduced complex}
The complete backend uses the Frobenius algebra
\[
 A=\F[x]/(x^2),\qquad \deg(1)=1,\quad\deg(x)=-1,
\]
with multiplication and comultiplication
\begin{align*}
 m(1\otimes1)&=1,&m(1\otimes x)&=m(x\otimes1)=x,&m(x\otimes x)&=0,\\
 \Delta(1)&=1\otimes x+x\otimes1,&\Delta(x)&=x\otimes x.
\end{align*}
These are the usual cube-of-resolutions maps \cite{khovanov,barnatan}.
For a state $s\in\{0,1\}^n$, the smoothing at $(a,b,c,d)$ is
\[
 0:\ (a,b),(c,d),\qquad 1:\ (a,d),(b,c).
\]
Union-find on semiarcs determines its $k(s)$ circles. The circle containing arc
1 is marked; it is required to carry label $x$. The other circles carry either
$1$ or $x$. Thus the reduced state space has dimension $2^{k(s)-1}$.
For the crossing-free circle this dimension is one.

The code changes one state bit from zero to one at each cube edge and computes
which two circles merge or which circle splits by their incident semiarcs.
It applies $m$ or $\Delta$, copying all unchanged labels. Requiring $x$ on the
marked circle gives a subcomplex: a split of a marked $x$ has both new labels
$x$, and a nonzero merge involving a marked $x$ is again labelled $x$.

Cube signs vanish over $\F$. The Frobenius identities make the two composites
around every square agree; their sum is zero in characteristic two. Hence
$d^2=0$. The optional runtime check composes the implemented differential on
every generator and verifies this cancellation directly.

For a labelling with $r$ copies of $x$, the package uses
\[
 h=|s|,\qquad q=h+k(s)-2r.
\]
The differential increases $h$ by one and preserves $q$. These are deliberately
\emph{unnormalized} gradings. Orientation-dependent global shifts are
unnecessary for total rank; the program does not advertise a normalized
Khovanov-polynomial interface.

\subsection{Exact linear algebra and the decision theorem}
For each $(h,q)$, a differential column is a nonnegative Python integer whose
set bits name nonzero rows. Gaussian elimination stores one column for each
highest set-bit pivot; elimination is exact XOR. Let $c_{h,q}$ be the dimension
of a chain group and $r_{h,q}$ the rank of its outgoing differential. Then
\[
 \beta_{h,q}=c_{h,q}-r_{h-1,q}-r_{h,q},\qquad
 \beta=\sum_{h,q}\beta_{h,q}=C-2\sum_{h,q}r_{h,q}.
\]
The implementation checks nonnegativity and agreement of these two total-rank
calculations. Its decision is $\beta=1$ for the unknot, $\beta>1$ for a
nontrivial knot.

\begin{theorem}[Mathematical correctness of the complete backend]
For a valid classical one-component diagram $D$,
\[
 D\text{ represents the unknot}
 \quad\Longleftrightarrow\quad
 \dim_{\F}\Kh(D;\F)=1.
\]
Consequently the specified full-complex computation decides unknotness.
\end{theorem}
\begin{proof}
Kronheimer--Mrowka prove that reduced Khovanov homology of a knot has rank one
over $\Q$ precisely for the unknot \cite{km}. The reduced integral complex is
free in each chain degree. The universal coefficient theorem implies
\[
 \dim_{\Q}\Kh(D;\Q)\le \dim_{\F}\Kh(D;\F).
\]
The rational reduced homology is nonzero: its graded Euler characteristic is
a shifted reduced Jones polynomial, which is nonzero since its normalized
value at one is one. Therefore dimension one over $\F$ forces rank one over
$\Q$, and the cited detector gives the unknot. Conversely the unknot's reduced
complex has one generator, so its reduced homology over $\F$ has dimension
one. The implemented complex is this standard complex, up to grading shifts
and the diagram convention, neither of which changes its total dimension.
\end{proof}

This correctness argument uses established topology. It is not a replacement
for the Kronheimer--Mrowka proof, nor a mechanically checked verification of
all Python instructions. The code does not infer triviality from the Jones
polynomial, determinant, or a table of named knots.

\subsection{The actual upper bound is exponential}
\begin{proposition}
The full-complex reference backend has bit complexity $2^{O(n)}$ and uses at
most $2^{O(n)}$ bits of storage. It explicitly enumerates $2^n$ resolutions.
\end{proposition}
\begin{proof}
Smoothing a crossing can increase the number of components by at most one.
Starting with one knot component gives $k(s)\le n+1$. Thus
\[
 C=\sum_{s\in\{0,1\}^n}2^{k(s)-1}\le 2^n2^n=4^n.
\]
Each generator has at most $2n$ differential terms. Resolutions, incidence
maps and label maps require polynomial work in $n$ per state or term. A
conservative bound for elimination is $O(C^3\operatorname{poly}(n))$ bit
operations: there are at most $C$ columns, each uses at most $C$ pivot
eliminations, and an XOR handles at most $C$ row bits. Stored pivots use at
most $O(C^2\operatorname{poly}(n))$ bits. These bounds are exponential in $n$.
The optional $d^2$ check performs at most polynomially many differential
compositions per generator and does not change the exponential upper bound.
\end{proof}

The explicit cube enumeration rules out advertising the full-complex routine
as quasi-polynomial. Polynomial-time filters do not constitute a proof that
this fallback will be avoided on all large inputs. Resource caps make it a
sound partial procedure, not a faster total recognizer.

\section{Compressed cohomology-dual normal surfaces}
\subsection{A concrete finite triangulation format}
For tetrahedron $t$ and its face opposite vertex $f$, a gluing is either a
boundary marker or a pair $(u,p)$, where $p$ permutes $\{0,1,2,3\}$ and carries
face $f$ to face $p(f)$ in tetrahedron $u$. Reciprocal pairings are mandatory.
Different faces of one tetrahedron may be identified; a face cannot be paired
to itself.

The validator tracks global vertices and signed global edges. Identifying an
edge to its reverse is rejected. It checks a consistent orientation on the
tetrahedra. Vertex links are constructed as glued triangles; the valid
edge-pairing conditions give connected surface links. A closed link must have
Euler characteristic two and a link with boundary must have characteristic
one. Orientability then identifies them as spheres or disks. Positive-genus
ideal links are not accepted: this is a finite compact-manifold format, not an
ideal-triangulation interface.

\subsection{Computing cohomology directions}
An integral edge cochain $z$ is closed when, on every triangular face,
\[
 z_{ab}+z_{bc}-z_{ac}=0.
\]
Global edge orientations are included in these equations. A maximal spanning
forest in the vertex-edge graph fixes the gauge by making $z$ vanish on its
tree edges. Every rational cohomology class has a unique representative in
this gauge: an exact cochain can successively cancel the values on tree edges,
and a remaining exact tree-zero cochain vanishes everywhere.

Exact rational row reduction computes the kernel of the face equations in
non-tree variables. Clearing denominators and dividing a coordinate gcd gives
integral vectors whose classes form a basis over $\Q$. The package does
\emph{not} assert that these vectors form a saturated $\Z$-basis. In a matrix
of size $O(t)$ with bounded integer entries, determinant bounds control the
bit length of the rational row-reduction data polynomially in $t$; this step
does not require iterating over the values of its output coordinates.

\subsection{The normal vector without disk expansion}
In a tetrahedron set $h_0=0$ and $h_i=z_{0i}$ for $i=1,2,3$. Closedness gives
$z_{ij}=h_j-h_i$. Regard the affine height function modulo $\Z$. Across a
paired face the two local height functions differ by an integer constant, so
their half-integral level sets agree.

Sort the four heights, breaking ties by vertex name. For $j=1,2,3$, let
\[
 d_j=h_{v_{j+1}}-h_{v_j}\ge0.
\]
Exactly $d_j$ half-integral levels separate the lowest $j$ vertices from the
rest. For $j=1$ they are triangles around $v_1$; for $j=3$ they are triangles
around $v_4$; and for $j=2$ they are quadrilaterals of a single type. The result
is a seven-entry vector in the order
\[
 (t_0,t_1,t_2,t_3,q_{01\mid23},q_{02\mid13},q_{03\mid12}).
\]
Only a constant number of integer comparisons and subtractions occurs per
tetrahedron. There is no loop over the $d_j$ individual disks. The three normal
arc counts on each paired face are checked explicitly; at most one
quadrilateral type per tetrahedron is allowed.

\begin{proposition}
For a closed integral cochain $z$ on a valid input triangulation, this
construction gives a properly embedded, cooriented normal surface dual to $z$,
possibly disconnected or empty. Its coordinates are computed in time
polynomial in the triangulation size and the bit length of $z$.
\end{proposition}
\begin{proof}
The map to $\mathbb R/\mathbb Z$ just described is well-defined because the
height changes across gluings are integral. The half-integral fiber avoids
vertices. On each tetrahedron it consists of parallel planar disks, of exactly
the types and multiplicities counted above. The normal matching check verifies
their gluing across faces. Increasing height supplies a consistent
coorientation; the integral intersection number with an oriented edge is its
cochain value. This is the cohomology-duality characterization of the fiber.
The sorting and integer operations have the stated bit cost.
\end{proof}

Euler characteristic is also obtained without expansion. If $F$ is the total
number of normal disks, $E$ the sum of normal arc counts over global faces
(counting a paired face once), and $V=\sum_e|z(e)|$, then
\[
 \chi(S)=V-E+F.
\]
These are the vertex, edge and face counts of the induced surface cell
structure, including boundary cells.

In the one-tetrahedron solid-torus example, the computed cocycle is
$(3,-1,-2)$, with normal vector $(1,1,0,0,0,0,1)$ and Euler characteristic one.
Scaling the cocycle by $10^{1000}$ produces $3\cdot10^{1000}$ normal disks but
still only seven integer coordinates. This is a compression demonstration,
not a genus bound or a full implementation of the Agol--Hass--Thurston methods
mentioned in the notes.

\subsection{The contract deliberately stops here}
A nonzero cohomology direction gives a nonzero relative homology class, but
this implementation does not extract a connected nonseparating component.
Different basis directions are not assembled into a controlled simultaneous
multisurface. The code does not cut the manifold, transport its boundary
pattern, retain its embedding in $S^3$, or establish the quadratic complexity
bound of the notes. A short binary normal vector by itself proves none of
those properties.

\section{Essential boundary patterns on a known ball}
The source definition asks whether every properly embedded disk meeting the
pattern at most three times has an allowed boundary disk: empty pattern,
one arc, or one tripod \cite[pp.~24--25]{slides}. The implementation handles
this terminal predicate when the ambient manifold is already known to be a
3-ball. It does not decide whether an arbitrary manifold is a ball.

A trivalent pattern is represented by three cyclically ordered darts per
vertex and a perfect matching of darts into edges. Loops and parallel primal
edges are permitted. The permutation of face walks determines the spherical
embedding, which is checked component by component. Isolated circle components
are counted separately.

An empty pattern and one simple closed curve are essential. A disconnected
nonempty pattern is inessential, regardless of its nesting on the sphere.
For a connected trivalent graph, construct the embedded dual graph. The
following exhaustive short-cycle test is implemented:
\begin{enumerate}
\item A dual loop supplies a violating one-edge cycle.
\item Two distinct dual edges with the same endpoints supply a violating
two-edge cycle.
\item Otherwise enumerate dual triangles. A triangle that is not the boundary
of a dual face on either side supplies a violating three-edge cycle.
\end{enumerate}
The output records the dual vertices and crossed primal edges. The two-face
triangulation with dual graph $K_3$ is handled correctly; its triangles are
facial, so it is not incorrectly rejected by an unqualified 4-connectivity
criterion. The alternative characterization is a simple 4-connected dual,
or $K_3$ or $K_4$ \cite[Proposition~2.6]{lackenby2026}.

Here is a direct justification of the implemented cycle criterion. For a
connected pattern, a disjoint curve bounds an empty-pattern disk on one side.
After dual loops and parallel edges are excluded, a curve crossing the pattern
twice either gives an already excluded pair of distinct dual edges or crosses
the same primal edge twice, isolating an arc. A length-three dual walk in a
simple graph must be a triangle. Its curve isolates a tripod exactly when the
triangle is facial on one side. Thus a nonfacial triangle is the remaining
obstruction. Conversely each recorded short cycle can be realized by a simple
curve in the indicated spherical faces and pushed into the ball to span a
disk. The test is polynomial; a conservative bound for triangle enumeration
is cubic in the explicitly represented graph size.

The independent test implementation instead checks graph simplicity and
vertex connectivity by enumerating vertex separators of size at most three.
Agreement between these differently structured predicates is tested on all
four-vertex dart matchings and on larger sampled spherical patterns.

A short dual cycle is not yet the compression disk in an earlier manifold.
Recovering that disk through all cuts, with its pattern history and compressed
representation, is a separate unimplemented operation. Also, $b_1(M)=0$ alone
does not establish that $M$ is a ball: for example $S^2\times[0,1]$ has that
Betti number but is not a ball. A full hierarchy implementation must maintain
the topological hypotheses behind its terminal branch, rather than silently
promote a homology test into ball recognition.

\section{What a genuine quasi-polynomial proof must control}
\subsection{The exact arithmetic part}
The slides' digit-counting argument is schematic. For an actual counter with
retained digits $0\le a_i\le g$, pad the hierarchy to a fixed length $L$ and use
\[
 B=g+1,\qquad \Phi(a)=\sum_{i=1}^L a_iB^{L-i},\qquad
 0\le\Phi(a)\le B^L-1.
\]
The distinction between $g$ and $g+1$ matters: base $g$ cannot accommodate a
digit equal to $g$ with the required strictness.

\begin{lemma}[Suffix rebuilding]
Suppose the prefix through index $j-1$ is unchanged, the digit at $j$ decreases
by at least one, and every later digit remains between zero and $g$. Then the
potential decreases by at least one, even if the entire suffix is rebuilt.
\end{lemma}
\begin{proof}
The loss at $j$ is at least $B^{L-j}$. The greatest possible gain in the suffix
is
\[
 g\sum_{r=0}^{L-j-1}B^r=B^{L-j}-1.
\]
Subtracting gives a net decrease of at least one. For $g=0$ no such strict
update can occur, so the assertion is vacuous.
\end{proof}

\code{HierarchyBudget} evaluates this potential and checks a proposed update.
Comparisons are between corresponding completed states, not every individual
extension of a partial hierarchy. This distinction avoids falsely demanding
that appending a new positive digit lower the potential.

\subsection{The conditional running-time statement}
\begin{proposition}[A sufficient quantitative contract]
Assume an implemented geometric algorithm has at most $R(n)$ phases, at most
$L(n)+1$ elementary calls between comparable completed states, and a strict
potential drop between successive nonterminating completed states of each
phase. Assume its digits are bounded by $g(n)$ and its elementary calls,
including manipulation of all encoded data, cost at most $P(n)$ bit operations.
Then its work is bounded by
\[
 R(n)P(n)(L(n)+1)(g(n)+1)^{L(n)}.
\]
In particular, polynomial $R,P,g$ and logarithmic $L$ imply the requested
quasi-polynomial bound.
\end{proposition}
\begin{proof}
There are at most $(g+1)^L$ completed states in a strict nonnegative descent
phase. Multiply this count by the stipulated work between them and by the
number of phases. If $g\le n^c$ and $L\le C\log(n+2)$, the logarithm of the
exponential factor is $O((\log(n+2))^2)$. The remaining polynomial factors do
not change this order.
\end{proof}

This proposition is a contract, not a proof that the missing topology
satisfies it. In particular, an integer returned by
\code{conditional_operation_bound} does not certify the supplied values of
$L,g,R,P$.

\subsection{The unresolved geometric and bit-complexity obligations}
\begin{longtable}{@{}p{0.24\linewidth}p{0.70\linewidth}@{}}
\toprule
Requirement & What an implementation and proof still need to provide\\
\midrule
\endhead
Initial encoded geometry &
A finite layered handle structure for the knot exterior, a generalized
Heegaard/Morse description, boundary-pattern data, and the ambient information
needed by the bounded-complexity surface construction. This archive has no
PD-to-layered-exterior pipeline.\\[4pt]
Hierarchical multisurfaces &
Simultaneous properly embedded oriented surfaces with independent classes,
controlled intersections, the required pattern-complexity bounds, and the
compressed combinatorial data needed by subsequent cuts. Separate cohomology
vectors are not enough.\\[4pt]
Compressed modifications &
Cutting, component selection, compression and pattern-compression without
expanding large normal coordinates, together with exact identification maps
back through the hierarchy and correct retention or discarding of later
surfaces.\\[4pt]
Cheeger regions and reduction &
A constructive implementation of the alternatives on pages 101--109 of the
notes, including modifications of the generalized Heegaard splitting,
Haken-lemma and weak-reduction steps, with a quantitative progress measure
for restarts.\\[4pt]
Combined bounds &
A logarithmic hierarchy-depth bound valid after every rebuild, bounded digit
complexity, bounded representation size, and a bound for all phases and
subroutine bit costs. None follows merely from naming the operations in a
flowchart.\\
\bottomrule
\end{longtable}
The three implemented hierarchy components address only parts of this table.
They are not treated as an oracle-based realization of the whole algorithm.

\section{Verification results and reproducibility}
The delivered version passed 64 \code{unittest} test methods, including many
subcases. The core tests check diagram rejection, exact determinant arithmetic
against the Leibniz formula, binary elimination against dense row elimination,
and an independently structured small-diagram homology implementation. The
latter uses graph walks instead of union-find, label tuples instead of masks,
and dense byte matrices instead of packed columns.

Topological consistency tests include crossing-preserving rotations, marked
arc changes, mirror examples, generator cancellation, braid conjugation, the
braid relation, and positive and negative Markov stabilization. They also
check $d^2=0$, both recognition outcomes with determinant one, and the
non-verdict semantics of resource limits. These are finite tests, not a formal
verification of all inputs.

For boundary patterns, all 10,395 perfect matchings of the 12 darts at four
trivalent vertices are generated; nonspherical rotation systems are rejected,
and accepted systems are checked against the independent connectivity
criterion. Larger randomized systems and explicit spherical bipyramids test
separating triangles. Normal-coordinate tests include all 24 relabellings of
the solid-torus tetrahedron, face matching, sign reversal, and huge compressed
coordinates.

Fifteen example diagrams were run through the default recognizer and also
through full homology with the $d^2$ check enabled. A selection of actual
outputs is shown below. Here $n$ is the crossing count of the supplied
presentation, not necessarily the crossing number of the knot.
\begin{center}
\begin{tabular}{@{}lrrrr@{}}
\toprule
Example & $n$ & Determinant & $\dim_{\F}\Kh$ & Generators\\
\midrule
Crossing-free circle & 0 & 1 & 1 & 1\\
Trefoil & 3 & 3 & 3 & 15\\
Figure-eight & 4 & 5 & 5 & 33\\
Conjugated unknot & 6 & 1 & 1 & 369\\
$T(3,4)$ & 8 & 3 & 5 & 801\\
$T(3,5)$ & 10 & 1 & 7 & 4209\\
Alternate $10_{124}$ braid & 10 & 1 & 7 & 14973\\
\bottomrule
\end{tabular}
\end{center}
The conjugated unknot is the closure of
$\sigma_2^{-2}(\sigma_1\sigma_2)\sigma_2^2$. Undoing the conjugation and two
Markov destabilizations proves that it is trivial independently of the
calculation. The descending test does not recognize this particular diagram,
so it exercises the positive fallback path. The $T(3,5)$ example exercises
the negative fallback path and prevents treating determinant one as a
complete test.

Six independent braid/determinant data pairs were obtained through Wolfram
\code{KnotData} \cite{wolfram}. The returned data and the expression that
regenerates them are stored in \code{results/}. The runtime has no Wolfram
dependency. In this environment the code itself was tested locally with
CPython 3.13; no Lean checking or remote Python execution is claimed.

The selected $T(3,5)$ presentation used 1024 cube states and 4209 generators;
full homology plus $d^2$ checking took approximately 0.078 seconds in the
recorded run. The alternate presentation had the same answer but 14973
generators and approximately 0.299 seconds. These timings are illustrative
measurements of very small inputs. They give no evidence that the general
algorithm meets the requested asymptotic target.

The files \code{results/tests.log}, \code{results/benchmark.json}, and
\code{results/benchmark.log} contain the actual test and measurement output.
Running \code{benchmark.py} regenerates the benchmark files, including
environment information. The mathematical specification above and
\code{STATUS.json} are explicit about what this delivery does and does not
establish.

\section{Conclusion}
The archive contains real executable recognition code and useful,
independently tested pieces of the hierarchy approach. The general recognizer
is exact but exponential, and the geometric pieces remain disconnected from
the bounded-depth construction. Accordingly, this delivery is not a solution
to the requested $n^{O(\log n)}$ implementation. The remaining work consists of
the stated constructions and quantitative proofs, not of renaming the
fallback or replacing those constructions by unimplemented calls.

\begin{thebibliography}{9}
\bibitem{slides}
Marc Lackenby, \emph{Unknot recognition in quasi-polynomial time}, February
2021, 109-page talk supplied with the request as
\code{quasipolynomial-talk.pdf}. Page references in this report refer to that
specific attachment, including its repeated build-up slides.

\bibitem{lackenby2026}
Marc Lackenby, \emph{Incompressible surfaces, hierarchies and unknot
recognition}, arXiv:2607.23350v1, 25 July 2026.
\url{https://arxiv.org/abs/2607.23350}.

\bibitem{khovanov}
Mikhail Khovanov, \emph{A categorification of the Jones polynomial}, Duke
Mathematical Journal \textbf{101} (2000), 359--426.
\url{https://arxiv.org/abs/math/9908171}.

\bibitem{barnatan}
Dror Bar-Natan, \emph{On Khovanov's categorification of the Jones polynomial},
Algebraic \& Geometric Topology \textbf{2} (2002), 337--370.
\url{https://arxiv.org/abs/math/0201043}.

\bibitem{km}
P. B. Kronheimer and T. S. Mrowka, \emph{Khovanov homology is an unknot-detector},
Publications Math\'ematiques de l'IH\'ES \textbf{113} (2011), 97--208.
\url{https://arxiv.org/abs/1005.4346}.

\bibitem{kolay}
Sudipta Kolay, \emph{Knot Colorings: Coloring and Goeritz matrices},
arXiv:1910.08044, 2019. \url{https://arxiv.org/abs/1910.08044}.

\bibitem{descending}
Maciej Mroczkowski, \emph{Diagrammatic unknotting of knots and links in the
projective space}, arXiv:math/0207302, 2002.
\url{https://arxiv.org/abs/math/0207302}.

\bibitem{wolfram}
Wolfram Research, \emph{KnotData}, Wolfram Language documentation and live
braid/determinant data retrieved for this implementation.
\url{https://reference.wolfram.com/language/ref/KnotData.html}.
\end{thebibliography}
\end{document}
